\documentclass{article}
\usepackage[utf8]{inputenc}
\usepackage{graphicx}
\usepackage{amsmath}
\usepackage{amssymb}
\usepackage{url}
\usepackage{tabularx}
\usepackage[colorlinks=true,allcolors=blue]{hyperref}
\usepackage{caption}
\usepackage{pdfpages}
\usepackage{multicol}
\usepackage[ngerman]{babel}
\usepackage{multicol}
\usepackage{geometry}	% [textwidth=17cm,textheight = 24.5cm]
% Um die Dokumentränder zu sehen: [showframe] vor {geometry} einfügen
\usepackage{bbold}
\usepackage{makecell}
\usepackage{float}
\usepackage{braket}
\usepackage{cancel}
\usepackage{braket}
\usepackage{mathtools}
\usepackage{xltabular}
\usepackage{float}
\usepackage{MnSymbol}
\usepackage{multirow}
\usepackage{chemformula}

\renewcommand{\cellalign}{tl}

\setlength{\columnsep}{1cm}
\setlength\parindent{0pt}

\newcommand{\diff}{\mathrm{d}}
\newcommand{\konst}{\mathrm{konst.}}
\newcommand{\lr}[1]{\left(#1\right)}
\newcommand{\lreck}[1]{\left[#1\right]}
\newcommand{\dkon}[3]{\left(\frac{\partial #1}{\partial #2}\right)_{#3}}

\definecolor{gruen}{HTML}{038133}

\newcommand{\commenttext}[1]{\textcolor{gruen}{\text{#1}}}
\newcommand{\comment}[1]{{\color{gruen} {#1}}}

\begin{document}

\begin{center}
	\Large{\textbf{Allgemeine Relativitätstheorie}}\\
	\small{613348}\\
	\small{\today}
\end{center}

\tableofcontents


\newpage


\section{Spezielle Relativitätstheorie}
\subsection{Lorentz-Boost}
Für einen Lorentz-Boost in $x$-Richtung, gilt
	\begin{align}
		t' &= \gamma \lr{ t - \frac{\beta}{c} x } \\
		x' &= \gamma (x - \beta c t) \\
		y' &= y \\
		z' &= z
	\end{align}

mit dem $\beta$ und Lorentz-Faktor
	\begin{align}
		\beta &= \frac{v}{c} \\
		\gamma &= \frac{1}{\sqrt{1 - \lr{ \frac{v}{c}}^2 } } \equiv \frac{1}{\sqrt{1 - \beta^2 }}
	\end{align}

In Matrixschreibweise ergibt sich für die Lorentz-Transformation und die dazu inverse Lorentz-Transformation \textcolor{gruen}{(die man durch $x'\leftrightarrow x, t'\leftrightarrow t, v\to -v$ erhält)})
	\begin{align}
		\begin{pmatrix}
		ct' \\ x'
		\end{pmatrix}
		&=
		\gamma
		\begin{pmatrix}
			1 & -\beta \\
			-\beta & 1
		\end{pmatrix}
		\begin{pmatrix}
		ct \\ x
		\end{pmatrix} \\
		\Leftrightarrow \begin{pmatrix}
		ct \\ x
		\end{pmatrix}
		&=
		\gamma
		\begin{pmatrix}
			1 & +\beta \\
			+\beta & 1
		\end{pmatrix}
		\begin{pmatrix}
		ct' \\ x'
		\end{pmatrix}
	\end{align}

Das $\chi$ in der Herleitung des Boost-Transformation wird als Rapidität bezeichnet und hat die Eigenschaft
	\begin{equation}
		\theta = \text{arctanh}\frac{v}{c} \longrightarrow
		\begin{cases}
			0, & v\to 0 \\
			\infty, & v\to c
		\end{cases}
	\end{equation}

\vspace{0.5cm}

{\color{violet}
\hrule
	\fbox{Herleitung}\\
	Das System $\Sigma'$ (mit den Koordinaten $t',x',y',z'$) bewege sich mit der Geschwindigkeit $v$ in $x$-Richtung relativ zum System $\Sigma$ ($t,x,y,z$).\\
	Nehme an, dass $y',z'$ vom Boost unverändert sind und $t',x'$ daher auch nicht von $y',z'$ abhängen.\\
	Nehme weiter an, dass der Zusammenhang zwischen $t',x'$ und $t,x$ linear ist. \textcolor{red}{Warum?}
		\begin{align*}
			\begin{pmatrix}
				ct'\\ x'
			\end{pmatrix}
			&=
			\gamma
			\begin{pmatrix}
				L_{00} & L_{01} \\
				L_{10} & L_{11}
			\end{pmatrix}
			\begin{pmatrix}
				ct\\ x
			\end{pmatrix}
		\end{align*}
	
	Benutze nun die Invarianz des Vierervektors
		\begin{align*}
			-(ct')^2 + (x')^2 &= -( L_{00} ct + L_{01} x )^2 + (L_{10} ct + L_{11} x)^2 = \\
			&= -(L_{00}^2 (ct)^2 + L_{01}^2 x^2 + 2L_{00} L_{01} ct x)	+ (L_{10}^2 (ct)^2 + L_{11}^2 x^2 + 2L_{10}L_{11} ct x) = \\
			&= -(ct)^2 + x^2
		\intertext{Führe einen Koeffizientenvergleich bzgl. \fbox{$(ct)^2$}, \fbox{$x^2$} und \fbox{$2ct x$}}
			-1 &= -L_{00}^2 + L_{10}^2 \\
			1 &= \phantom{-}L_{01}^2 + L_{11}^2 \\
			0 &= -L_{00}L_{01} + L_{10}L_{11}
		\end{align*}

	Dieses Gleichungssystem wird gelöst durch $L_{00} = L_{11} = \cosh\chi$ und $L_{10} = L_{01} = \sinh\chi$. \textcolor{gruen}{Kann man sich merken, weil die Matrix der Lorentz-Transformation die Form einer Drehmatrix hat, nur mit sinh und cosh statt sin und cos.}
		\begin{align*}
			\begin{pmatrix}
				ct'\\ x'
			\end{pmatrix}
			&=
			\begin{pmatrix}
				\cosh\chi & -\sinh\chi \\
				-\sinh\chi & \cosh\chi
			\end{pmatrix}
			\begin{pmatrix}
				ct\\ x
			\end{pmatrix}
		\end{align*}
	
	Betrachte nun den Ursprung des Systems $\Sigma'$, d.h. $x'=0$. Dieser Punkt wird im System $\Sigma$ durch $x=vt$ beschrieben. Löse dieses Gleichungssystem:
		\begin{align*}
			x' &= 0 \tag{1}\\
			x &= v\cdot t \stackrel{\comment{1}}{=} \frac{v}{c}\cdot (ct) \tag{2}\\
		\intertext{Setze für $x'$ die Matrixmultiplikation ein und verwende dann Gl. (2), um $x$ durch $v$ auszudrücken.}
			\Rightarrow 0 &= x' = (-\sinh\chi \cdot ct + \cosh\chi \cdot x) = ct \underbrace{ \big( \cosh\chi \cdot \underbrace{\frac{x}{ct}}_{\stackrel{(2)}{=} \frac{v}{c}} - \sinh\chi \big) }_{\stackrel{!}{=}0}
		\intertext{Die Forderung, dass hintere Faktor null sein muss kommt daher, dass im allgemeinen $t\neq 0$ gilt.}
			\Rightarrow 0 &= \cosh\chi \cdot \frac{v}{c} - \sinh\chi \\
			\Leftrightarrow \frac{\sinh\chi}{\cosh\chi} = \tanh\chi &= \frac{v}{c}
		\end{align*}
	
	Nutze schließlich trigonometrische Funktionen, um $\sinh$ und $\cosh$ auszudrücken.
		\begin{align*}
			\cosh\chi &= \frac{1}{\sqrt{1 - \tanh^2\chi}} = \frac{1}{\sqrt{1 - \lr{ \frac{v}{c} }^2 }} \equiv \gamma \\
			\sinh\chi &= \tanh\chi \cdot \cosh\chi = \frac{v}{c} \cdot \frac{1}{\sqrt{1 - \lr{ \frac{v}{c} }^2 }} \equiv \beta\gamma
		\end{align*}
	
	Somit gilt für den Lorentz-Boost
		$$\begin{pmatrix}
		ct' \\ x'
		\end{pmatrix}
		=
		\gamma
		\begin{pmatrix}
			1 & -\beta \\
			-\beta & 1
		\end{pmatrix}
		\begin{pmatrix}
		ct \\ x
		\end{pmatrix}
		$$
	
\hrule
}


\newpage


\subsection{Viererimpuls und relativistische Energie}
Der Viererimpuls ist definiert als
	\begin{equation}
		p = \lr{ \frac{E}{c}, \vec{p} }
	\end{equation}

Aus dem Quadrat des Viererimpulses ergibt sich die relativistische Energie
	\begin{equation}
		E = \sqrt{(m c^2)^2 + (c \vec{p})^2 }
	\end{equation}

\vspace{0.5cm}

{\color{violet}
\hrule
	\fbox{Herleitung}\\
	Das Quadrat des Viererimpulses ist eine Invariante, d.h. sie ist in allen Bezugssystemen gleich. Bestimme sie daher in einem einfach auszurechnenden System: Dem Ruhesystem.\\
	Dort hat das Teilchen keine Geschwindigkeit und somit auch keinen Impuls.
		\begin{align*}
			p^2 &= -\frac{E^2}{c^2} + 0
		\intertext{Die Energie entspricht der Ruhenergie des Teilchens $E=mc^2$.}
			\Rightarrow p^2 &= -\frac{(mc^2)^2}{c^2} = -m^2c^2
		\end{align*}

	Setze dieses $p^2 = -m^2c^2$ in die allgemeine Formel für $p^2$ ein
		\begin{align*}
			p^2 &= -\frac{E^2}{c^2} + \vec{p}^2 \stackrel{!}{=} -m^2c^2 \\
			\Leftrightarrow \frac{E^2}{c^2} &= \vec{p}^2 + m^2c^2 \\
			\Leftrightarrow E^2 &= (c\vec{p})^2 + (mc^2)^2
		\end{align*}


\hrule
}




\newpage


\section{Grundlegende Formeln und Konventionen in der ART}
Gemäß der Einstein'schen Summenkonvention wird über gleichnamige Indizes, die paarweise oben bzw. unten stehen summiert
	\begin{align}
		A_{...,\mu,...} B^{...,\mu,...} &= \sum_{\mu=0}^3 A_{...,\mu,...} B^{...,\mu,...} \\
		A^{...,\mu,...} B_{...,\mu,...} &= \sum_{\mu=0}^3 A^{...,\mu,...} B_{...,\mu,...}
	\end{align}

Kontravariante und kovariante Vektoren können wie folgt über die Metrik ineinander umgerechnet werden
	\begin{align}
		x^\mu &= g^{\mu\nu} x_\nu \\
		x_\mu &= g_{\mu\nu} x^\nu
	\end{align}

Die Minkowski-Metrik im flachen Raum und für kartesische Koordinaten lautet
	\begin{equation}
		g_{\mu\nu} = 
		\begin{pmatrix}
			-1 & 0 & 0 & 0 \\
			0 & 1 & 0 & 0 \\
			0 & 0 & 1 & 0 \\
			0 & 0 & 0 & 1
		\end{pmatrix}
	\end{equation}

Das Kronecker-Delta lautet
	\begin{equation}
		g^{\mu{\color{gruen}\lambda}} g_{{\color{gruen}\lambda}\nu} = \delta^\mu_\nu
	\end{equation}

Das invariante Längenelement im Minkowski-Raum ist
	\begin{equation}
		\diff s^2 := {\color{gruen} \diff x_\mu \diff x^\mu} = -(\diff x^0)^2 + (\diff \vec{x})^2
	\end{equation}

Die Eigenzeit eines Objekts ist
	\begin{equation}
		\diff \tau = \sqrt{-\diff s^2}
	\end{equation}


\newpage


\section{Koordinatentransformationen}
Die Koordinatentransformation $x \leftrightarrow x'$ werden durch die Transformationsmatrix $S$ durchgeführt und lauten
	\begin{alignat}{2}
		x'^\mu &:= \frac{\partial x'^\mu}{\partial x^\alpha} \cdot x^\alpha &&\equiv  S^\mu{}_\alpha \cdot x^\alpha\\
		x^\mu &:= \frac{\partial x^\mu}{\partial x'^\alpha} \cdot x'^\alpha &&\equiv (S^{-1})^\mu{}_\alpha \cdot x'^\alpha
	\end{alignat}

\textcolor{gruen}{Merke: Die Transformationen entstehen durch "Hinzufügen einer Eins", d.h. Kettenregel. Und $S$ ist definiert als die Transformation von $x$ nach $x'$.}

Die Metrik transformiert gemäß
	\begin{align}
		g'_{\mu\nu} &= g_{\rho\sigma} \cdot \frac{\partial x^\rho}{\partial x'^\mu} \frac{\partial x^\sigma}{\partial x'^\nu} \\
		g'^{\mu\nu} &= g^{\rho\sigma} \cdot \frac{\partial x'^\mu}{\partial x^\rho} \frac{\partial x'^\nu}{\partial x^\sigma}
	\end{align}

{\color{violet}
\hrule
	\fbox{Herleitung}\\
	Die Transformation der Metrik ergibt sich aus der Invarianz des Linienelements $\diff s^2$ unter Koordinatentransformation.
		\begin{align*}
			ds^2 = g'_{\mu\nu} \cdot \diff x'^\mu \cdot \diff x'^\nu &= g_{\rho\sigma} \cdot \diff x^\rho \cdot \diff x^\sigma = g_{\rho\sigma} \cdot \lr{ \frac{\partial x^\rho}{\partial x'^\alpha} \diff x'^\alpha } \cdot \lr{ \frac{\partial x^\sigma}{\partial x'^\beta} \diff x'^\beta } \\
			&\stackrel{\comment{\alpha\to\mu,\beta\to\nu}}{=} g_{\rho\sigma} \cdot \frac{\partial x^\rho}{\partial x'^\mu} \frac{\partial x^\sigma}{\partial x'^\nu} \cdot \diff x'^\mu \diff x'^\nu \\
			\Rightarrow g'_{\mu\nu} &= g_{\rho\sigma} \cdot \frac{\partial x^\rho}{\partial x'^\mu}\frac{\partial x^\sigma}{\partial x'^\nu}
		\end{align*}
\hrule
}

\subsection{Tensoren}
Ein (m,n)-Tensor \textcolor{gruen}{($m$ Indizes oben und $n$ Indizes unten)} ist ein Objekt, das wie folgt transformiert:
	\begin{equation}
		T'^{\mu_1,...,\mu_m}{}_{\nu_1,...,\nu_n} = S^{\mu_1}{}_{\alpha_1} \cdot ... \cdot S^{\mu_m}{}_{\alpha_m} \cdot T^{\alpha_1,...,\alpha_m}{}_{\beta_1,...,\beta_n} \cdot (S^{-1})^{\beta_1}{}_{\nu_1} \cdot ... \cdot (S^{-1})^{\beta_m}{}_{\nu_m}
	\end{equation}

\subsubsection{Kontravariante und kovariante Vektoren}
Ein kontravarianter Vektor $W^\mu$ bzw. kovarianter Vektor $W_\mu$ ist ein (1,0)- bzw. (0,1)-Tensor und somit definiert über die Transformation
	\begin{alignat}{2}
		W'^\mu(x') &= \frac{\partial x'^\mu}{\partial x^\alpha} \cdot W^\alpha(x) &&= S^\mu{}_\alpha \cdot W^\alpha(x) \\
		W'_\mu(x') &= \frac{\partial x^\alpha}{\partial x'^\mu} \cdot W_\alpha(x) &&= {(S^{-1})}^\alpha{}_\mu \cdot W_\alpha(x)
	\end{alignat}

\subsubsection{Skalar}
Ein Skalar ist ein (0,0)-Tensor und somit definiert über
	\begin{equation}
		\phi'(x') = \phi(x)
	\end{equation}


\newpage


\section{Linienelement und Metrik}
Das Linienelement $\diff s$, das in der Raum-Zeit zurückgelegt wird, ist durch die Metrik $g_{\mu\nu}$ mit den infinitesimalen $\diff t, \diff x, \diff y, \diff z$, die die Bewegung in der Raumzeit beschreiben, verknüpft.
	\begin{equation}
		ds^2 = g_{\mu\nu} \diff x^\mu \diff x^\nu
	\end{equation}

Die Metrik hängt dabei von der Raum-Zeit, sowei den Koordinaten ab, in denen man sie betrachtet.

	\begin{itemize}
		\item $\mathbb{R}^2$ kartesische Koordinaten $(x,y)$
			$$g =
			\begin{pmatrix}
				1 & 0 \\
				0 & 1
			\end{pmatrix}$$
		
		\item $\mathbb{R}^2$ Polarkoordinaten $(r,\phi)$
			$$g =
			\begin{pmatrix}
				1 & 0 \\
				0 & r^2
			\end{pmatrix}$$
		
		\item $\mathbb{R}^3$ Kugelkoordinaten $(r,\theta,\varphi)$
			$$g =
			\begin{pmatrix}
			1 & 0 & 0 \\
			0 & r^2 & 0 \\
			0 & 0 & r^2\sin^2\theta
			\end{pmatrix}
			$$
		
		\item $\mathbb{R}^{1,3}$ in Kugelkoordainten $(t,r,\theta,\varphi)$\\
		\textbf{De-Sitter-Metrik}
			$$g = 
			\begin{pmatrix}
				-1 & 0 & 0 & 0 \\
				0 & a^2(t) & 0 & 0 \\
				0 & 0 & a^2(t) & 0 \\
				0 & 0 & 0 & a^2(t)
			\end{pmatrix}
			$$
		
		mit $a(t) = e^{H t}$.
		
		\item $\mathbb{R}^{1,3}$ in Kugelkoordinate $(t,r,\theta,\varphi)$\\
		\textbf{Schwarzschild-Metrik}
			$$g = 
			\begin{pmatrix}
				-\lr{ 1 - \frac{2M}{r} } & 0 & 0 & 0 \\
				0 & \frac{1}{1 - \frac{2M}{r}} & 0 & 0 \\
				0 & 0 & r^2 & 0 \\
				0 & 0 & 0 & r^2\sin^2\theta
			\end{pmatrix}
			$$
		
		\textbf{Friedmann-Robertson-Walker-Metrik (FRW)}
			$$
			g_{\mu\nu} = 
			\begin{pmatrix}
				-1 & 0 & 0 & 0 \\
				0 & a(t)^2 \frac{1}{1 - k \frac{r^2}{L^2}} & 0 & 0 \\
				0 & 0 & a(t)^2r^2 & 0 \\
				0 & 0 & 0 & a(t)^2 r^2\sin^2\theta
			\end{pmatrix}
			$$
	\end{itemize}


\newpage


{\color{violet}
\hrule
	\fbox{Herleitung}\\
	ges: Metrik für $\mathbb{R}^2$ in Polarkoordinaten
		\begin{align*}
			x &= r\cos\varphi \\
			y &= r\sin\varphi
		\end{align*}
	
	Betrachte den $\mathbb{R}^3$ in kartesischen Koordinaten.
		$$ds^2 = \diff x^2 + \diff y^2$$

	Formuliere zunächst die Infinitesimale $dx$, $dy$, $dz$ in Kugelkoordinaten und setze sie anschließend in $ds^2$ ein. Fasse dafür $x,y,z$ als Funktionen von $r,\theta,\varphi$ auf. Kleine Änderungen sind dann gemäß Kettenregel
		\begin{align*}
			dx &= \frac{\partial x}{\partial r} \diff r + \frac{\partial x}{\partial\phi} \diff\phi = \cos\phi \cdot \diff r - r\sin\phi \cdot \diff\phi \\
			dy &= \frac{\partial y}{\partial r} \diff r + \frac{\partial y}{\partial\phi} \diff\phi = \sin\phi \cdot \diff r + r\cos\phi \cdot \diff\phi \\
			\Rightarrow ds^2 &= \diff x^2 + \diff y^2 =
			\lr{ \cos\phi \cdot \diff r - r\sin\phi \cdot \diff\phi }^2 + \lr{ \sin\phi \cdot \diff r + r\cos\phi \cdot \diff\phi }^2 = \\
			&= \cos^2\phi \diff r^2 + r^2\sin^2\phi \diff\phi^2 - \cancel{2r\cos\phi\sin\phi \diff r\diff \phi} + \\
			&+ \sin^2\phi \diff r^2 + r^2\cos^2\phi \diff\phi^2 + \cancel{2 r\sin\phi\cos\phi \diff r\diff \phi} = \\
			&= \lr{\cos^2\phi + \sin^2\phi}\diff r^2 + r^2\lr{\sin^2\phi + \cos^2\phi} \diff \phi^2 = \diff r^2 + r^2 \diff \phi^2
		\end{align*}
	
	Dies muss gleich sein mit
		\begin{align*}
			ds^2 &= g'_{\mu\nu} \diff x^\mu \diff^\nu = g_{r\phi} \diff x^r \diff x^\phi +  g_{\phi r} \diff x^r \diff x^\phi + g_{rr} (\diff x^r)^2 + g_{\phi\phi} (\diff x^\phi)^2 \stackrel{!}{=} \\
			&\stackrel{!}{=} \diff r^2 + r^2\diff\phi^2 \\
			\Rightarrow g_{r\phi} = g_{\phi r} &= 0\\
			\Rightarrow g_{rr} &= 1 \\
			\Rightarrow g_{\phi\phi} &= r^2\\
			\Rightarrow g_{\mu\nu} &= 
			\begin{pmatrix}
				1 & 0 \\
				0 & r^2
			\end{pmatrix}
		\end{align*}
\hrule
}


\newpage


\section{Intrinsische und extrinsische Krümmung}
	\begin{enumerate}
		\item Aufstellen der  Gl. für die Mannigfaltigkeit
		
		\item Verschieben des interessierenden Punktes in den Ursprung
		
		\item Umstellen nach einer Variable
		
		\item Taylor-Entwicklung bis 2. Ordnung
		
		\item Extrinsische, intrinsische Krümmmung aus Hesse-Matrix bestimmen
	\end{enumerate}


	\begin{itemize}
		\item \textbf{Intrinsische Krümmung}\\
		Die intrinsische Krümmung ist die Krümmung, die durch Messungen auf der Mannigfaltigkeit bestimmt werden kann.
			\begin{equation}
				K_{\text{int}} = \det M
			\end{equation}
		
		\item \textbf{Extrinsische Krümmung}\\
		Die extrinsische Krümmung beschreibt, wie die Mannigfaltigkeit in der nächsthöheren Dimension eingebettet ist.
			\begin{equation}
				K_{\text{ext}} = \frac{1}{4} \text{Sp}(M)^2
			\end{equation}
	\end{itemize}

\vspace{0.5cm}

{\color{blue}
\hrule
	\fbox{Beispiel}\\
	geg: Kugel im $\mathbb{R}^3$\\
	\\
	Für den Kugeloberfläche gilt
		$$x^2+y^2+z^2 = L^2$$
	
	Verschiebe die Kugel so, dass der interessierende Punkt, an dem die Krümmung bestimmt werden soll im Ursprung liegt. Verschiebe hier die Kugel entlang der $z$-Achse.
		$$x^2+y^2+(z-L)^2 = L^2$$
	
	\textcolor{gruen}{Für $z=0$ und $z=2L$ müssen $x=y=0$ sein. Die verschobene Kugel liegt also zwischen $0$ und $2L$ entlang der $z$-Achse.}
		\begin{figure}[H]
			\centering
			\includegraphics[width=0.4\textwidth]{Grafiken/ART Krümmung Kugel.pdf}
			\caption{Tangentialebene an einer Kugel}
		\end{figure}
		
	Stelle nun nach $z = z(x,y)$ um. Dabei muss sich überlegt werden, welches Vorzeichen sinnvoll ist.
		\begin{align*}
			z-L &= \pm\sqrt{L^2 - x^2 - y^2} \\
			\Leftrightarrow z &= L\pm\sqrt{L^2 - x^2 - y^2}
		\end{align*}
	
	Betrachte die Lösung, die die Krümmung nach oben, also den unteren Punkt der Kugel beschreibt.\\
	\\
	Mache nun eine Taylor-Entwicklung von $z$ um den Punkt $(x_0,y_0)=(0,0)$
		\begin{align*}
			z(x,y) &\approx \left.z(x,y)\right|_{0,0} + \lreck{ \left.\frac{\partial z(x,y)}{\partial x}\right|_{0,0} (x-0) + \left.\frac{\partial z(x,y)}{\partial y}\right|_{0,0} (y-0) } + \\
			 &+ \frac{1}{2} (x,y)
			 \left.
			\begin{pmatrix}
				\partial_x^2 z(x,y) & \partial_x\partial_y z(x,y) \\
				\partial_x\partial_y z(x,y) & \partial_y^2 z(x,y)
			\end{pmatrix}
			\right|_{0,0}
			\begin{pmatrix}
				x\\y
			\end{pmatrix} + ...
		\end{align*}

	Berechne die Ableitugen
		\begin{align*}
			\partial_x z &= \partial_x L - \lr{L^2 - x^2 - y^2}^{1/2} = -\frac{1}{2} \frac{1}{\sqrt{L^2 - x^2 - y^2}} (-2x) = \frac{x}{\sqrt{L^2 -x^2 - y^2}} \\
			\partial_x^2 z &= \partial_x \frac{x}{\sqrt{L^2 -x^2 - y^2}} = \frac{1}{\sqrt{L^2 -x^2 - y^2}} - \frac{2x^2}{\sqrt{L^2 -x^2 - y^2}}
		\end{align*}

	Einsetzen in die Taylor-Entwicklung liefert
		\begin{align*}
			z &\approx 0 + 0 + \frac{1}{2} (x,y)
			\begin{pmatrix}
				1/L & 0\\
				0 & 1/L
			\end{pmatrix}
			\begin{pmatrix}
				x\\y
			\end{pmatrix}
			\equiv \frac{1}{2} x^T M x
		\end{align*}

	Die intrinsische Krümmung ist somit
		$$K_{\text{int}} = \det M =
		\begin{vmatrix}
			1/L & 0\\
			0 & 1/L
		\end{vmatrix} = \frac{1}{L^2}$$
	
	Die extrinsische Krümmung ist
		$$K_{\text{ext}} = \frac{1}{4}\text{Sp} (M)^2 = \frac{1}{4} \lr{\frac{2}{L}}^2 = \frac{1}{L^2 }$$


\hrule
}


\newpage


\section{Geodätengleichung}
Die Geodätengleichung liefert die Bewegungsgleichungen $X^\rho(l)$, für ein Teilchen, das sich durch die durch $g_{\mu\nu}$ beschriebene gekrümmte Raumzeit bewegt.
	\begin{align}
		\frac{\diff^2X^\rho}{\diff l^2} + \Gamma^\rho_{\mu\nu} \frac{\diff X^\mu}{\diff l}\frac{\diff X^\nu}{\diff l} &= 0 \\
		\Leftrightarrow \frac{\diff V^\rho}{\diff l} + \Gamma^\rho_{\mu\nu} V^\mu V^\nu &= 0
	\end{align}

mit dem Christoffel-Symbol
	\begin{equation}
		\Gamma^\rho_{{\color{red}\mu}{\color{blue}\nu}} := \frac{1}{2}g^{\rho{\color{gruen}\sigma}} \lr{ \partial_{\color{red}\mu} g_{{\color{gruen}\sigma}{\color{blue}\nu}} + \partial_{\color{blue}\nu} g_{{\color{red}\mu}{\color{gruen}\sigma}} - \partial_{\color{gruen}\sigma} g_{{\color{red}\mu}{\color{blue}\nu}} }
	\end{equation}

Diese haben
	\begin{equation}
		D\cdot \frac{D(D+1)}{2}
	\end{equation}

unabhängige Komponenten.\\
\\
Für die Bestimmung der Christoffel-Symbole gibt es die Methoden
	\begin{itemize}
		\item Nutzen der Definition
		
		\item $D_\lambda g_{\mu\nu} = 0$\\
		Durch Einsetzen der kovarianten Ableitung folgt
			$$\partial_\lambda g_{\mu\nu} - \Gamma^\alpha_{\lambda\mu}g_{\alpha\nu} - \Gamma^\alpha_{\lambda\nu}g_{\mu\alpha} = 0$$
		
		\item Geodätengleichung\\
		Durch Einsetzen einer bekannten Lagrange-Funktion in die Euler-Lagrange-Gleichung und Vergleichen dieser DGLs mit denen aus der Geodätengleichung kann man die Christoffel-Symbole erkennen.
	\end{itemize}


\textbf{Staudachers ''Bonusgleichung''}
	\begin{equation}
		\mathcal{L} = 1 \qquad\Leftrightarrow\qquad g_{\mu\nu} \frac{\diff X^\mu}{\diff l}\frac{\diff X^\nu}{\diff l} = 1
	\end{equation}

\textcolor{gruen}{Z.B. für kartesische Koordinaten und flache Metrik: $\frac{d^2 x}{\diff l^2} + \frac{\diff^2 y}{\diff l^2} = 1$.}


\newpage


\vspace{0.5cm}

{\color{blue}
\hrule
	\fbox{Beispiel}\\
	ges: Geodätengleichungen für Bewegung auf 2d Kugel $S^2$ \textcolor{gruen}{(PÜ 5.2 a), b))}\\
	\\
	Die Metrik der 2d Kugel $S^2$ ist
		\begin{equation*}
			g_{\mu\nu} = 
			\begin{pmatrix}
				1 & 0 \\
				0 & \sin^2\theta
			\end{pmatrix}
		\end{equation*}

	Bestimme nun die nicht verschwindenden Christoffel-Symbole
		\begin{align*}
			\Gamma^\rho_{\mu\nu} &= \frac{1}{2}g^{\rho\sigma} \lr{ \partial_\mu g_{\sigma\nu} + \partial_\nu g_{\mu\sigma} - \partial_\sigma g_{\mu\nu} } \\
			\Rightarrow \Gamma^\theta_{\varphi\varphi} &= -\sin\theta\cos\theta \\
			\Rightarrow \Gamma^\varphi_{\theta\varphi} = \Gamma^\varphi_{\varphi\theta} &= \cot\theta
		\end{align*}

	Führe nun die Summationen in der Geodätengleichung aus und setze die gefundenen Christoffel-Symbole ein.
		\begin{align*}
			\frac{\diff^2 X^\theta}{\diff l^2} + \Gamma^\theta_{\mu\nu} \frac{\diff X^\mu}{\diff l}\frac{\diff X^\nu}{\diff l} = 0 \\
			\Leftrightarrow \frac{\diff^2 \theta}{\diff l^2} + \Gamma^\theta_{\varphi\varphi} \frac{\diff X^\varphi}{\diff l}\frac{\diff X^\varphi}{\diff l} = 0 \\
			\Leftrightarrow \frac{\diff^2\theta}{\diff l^2} - \sin\theta\cos\theta \lr{\frac{\diff\varphi}{\diff l}}^2 = 0
		\intertext{und}
			\frac{\diff^2 X^\varphi}{\diff l^2} + \Gamma^\varphi_{\theta\varphi} \frac{\diff X^\theta}{\diff l} \frac{\diff X^\mu}{\diff l} + \Gamma^\varphi_{\varphi\theta} \frac{\diff X^\theta}{\diff l} \frac{\diff X^\mu}{\diff l} &= 0 \\
			\frac{\diff \varphi^2}{\diff l^2} + 2\cot\theta \frac{\diff \theta}{\diff l} \frac{\diff \varphi}{\diff l} &= 0
		\end{align*}
\hrule
}


\newpage


{\color{violet}
\hrule
	\fbox{Herleitung}\\
	Die Geodätengleichung bestimmt den Weg zwischen einem Startpunkt A und einem Endpunkt B, für den das infinitesimale Linienelement in der Raum-Zeit $\diff s$ minimal wird.
		\begin{align*}
			\text{Länge}(A,B) &= \int \sqrt{ g_{\mu\nu} \diff X^\mu \diff X^\nu } = \int_{\lambda_A}^{\lambda_B} \diff\lambda\, \sqrt{ g_{\mu\nu} \frac{\diff X^\mu}{\diff\lambda} \frac{\diff X^\nu}{\diff\lambda} } \equiv \int_{\lambda_A}^{\lambda_B} \diff\lambda\, \sqrt{g_{\mu\nu} X'^\mu X'^\nu} \\
			&\equiv \int_{\lambda_A}^{\lambda_B} \diff\lambda\, \mathcal{L}
		\intertext{Dies ist äquivalent zur Euler-Lagrange-Gleichung}
			\Leftrightarrow 0 &= \frac{\diff}{\diff\lambda} \frac{\partial \mathcal{L}}{\partial X'^\sigma } - \frac{\partial \mathcal{L}}{\partial X^\sigma} \\
			\stackrel{\comment{\mathcal{L}}}{\Leftrightarrow} 0 &= \frac{\diff}{\diff\lambda} \frac{\partial \sqrt{g_{\mu\nu} X'^\mu X'^\nu} }{\partial X'^\sigma } - \frac{\partial \sqrt{g_{\mu\nu} X'^\mu X'^\nu} }{\partial X^\sigma} \\
			\Leftrightarrow 0 &= \frac{\diff}{\diff\lambda} \frac{1}{2\mathcal{L}} \frac{\partial (g_{\mu\nu} X'^\mu X'^\nu)}{\partial X'^\sigma} - \frac{1}{2\mathcal{L}} \frac{\partial (g_{\mu\nu} X'^\mu X'^\nu)}{\partial X^\sigma} \\
			0 &= \frac{\diff}{\diff\lambda} \frac{1}{2\mathcal{L}} \lr{ g_{\mu\nu} \delta^\mu_\sigma X'^\nu + g_{\mu\nu} X'^\mu \delta^\nu_\sigma } - \frac{1}{2\mathcal{L}} \frac{\partial (g_{\mu\nu})}{\partial X^\sigma} X'^\mu X'^\nu \\
			0 &= \frac{\diff}{\diff\lambda} \frac{1}{\cancel{2}\mathcal{L}} \cancel{2} g_{\sigma\nu} X'^\nu - \frac{1}{2\mathcal{L}} \partial_\sigma (g_{\mu\nu}) X'^\mu X'^\nu \\
		\intertext{Verwende nun $\lambda = l$ und somit $\mathcal{L}=1$.}
			0 &= \frac{\diff}{\diff l} (g_{\sigma\nu} X'^\nu) - \frac{1}{2} \partial_\sigma (g_{\mu\nu}) X'^\mu X'^\nu \\
			0 &= \underbrace{\frac{\diff g_{\sigma\nu}}{\diff l}}_{\comment{= \frac{\diff g_{\sigma\nu}}{\diff X^\mu} \frac{\diff X^\mu}{\diff l}}} X'^\nu + g_{\sigma\nu} \frac{\diff^2 X^\nu}{\diff l^2} - \frac{1}{2} \partial_\sigma (g_{\mu\nu}) X'^\mu X'^\nu \\
			0 &= \partial_\mu (g_{\sigma\nu}) X'^\mu X'^\nu + g_{\sigma\nu} \frac{\diff^2 X^\nu}{\diff l^2} - \frac{1}{2} \partial_\sigma (g_{\mu\nu}) X'^\mu X'^\nu \\
		\stackrel{\comment{\cdot g^{\rho\sigma}}}{\Leftrightarrow} 0 &= \delta^\rho_\nu \frac{\diff^2 X^\nu}{\diff l^2} + g^{\rho\sigma} \partial_\mu (g_{\sigma\nu}) X'^\mu X'^\nu - \frac{1}{2} g^{\rho\sigma} \partial_\sigma (g_{\mu\nu}) X'^\mu X'^\nu \\
		\Leftrightarrow 0 &= \frac{\diff^2 X^\rho}{\diff l^2} + \frac{1}{2} g^{\rho\sigma} \lr{ 2 \partial_\mu g_{\sigma\nu} - \partial_\sigma g_{\mu\nu} } X'^\mu X'^\nu
		\intertext{Nun kann man $2 \partial_\mu g_{\sigma\nu} = \partial_\mu g_{\sigma\nu} + \partial_\mu g_{\sigma\nu}$ schreiben und wegen der Kontraktion mit $X'^\mu X'^\nu$ in einem Term $\mu\leftrightarrow\nu$ tauschen. Dann lässt sich das Christoffel-Symbol erkennen.}
			\Leftrightarrow 0 &= \frac{\diff^2 X^\rho}{\diff l^2} + \frac{1}{2} g^{\rho\sigma} \lr{\partial_\mu g_{\sigma\nu} + \partial_\nu g_{\sigma\mu} - \partial_\sigma g_{\mu\nu} } X'^\mu X'^\nu \\
			\Leftrightarrow 0 &= \frac{\diff^2 X^\rho}{\diff l^2} + \Gamma^\rho_{\mu\nu} X'^\mu X'^\nu \qquad\blacksquare
		\end{align*}
\hrule
}


\newpage


{\color{violet}
\hrule
	\fbox{Herleitung}\\
	Zeige, dass das Prinzip der kleinsten Wirkung äquivalent zu den Euler-Lagrange-Gleichungen ist.\\
	\\
	Betrachte die Wirkung
		$$S[x^\mu,x'^\mu] = \int \diff l\, \mathcal{L}(t,x^\mu,x'^\mu)$$

	Angenommen $x^\mu$ sind Lösungen, die $S$ minimieren. Betrachte nun die Addition einer anderen Funktion $\eta^\mu$, multipliziert mit einem Faktor $\varepsilon$. Sie soll die Randbedingungen erfüllen $\eta^\mu(l_i) = \eta^\mu(l_f) = 0$, also, dass die Variation am Rand verschwindet und die variierte Funktion mit der unvariierten übereinstimmt.
		\begin{align*}
			x^\mu &\to x^\mu + \varepsilon \eta^\mu \\
			x'^\mu &\to x'^\mu + \varepsilon \eta'^\mu \\
			\Rightarrow S[x^\mu + \varepsilon \eta^\mu, x'^\mu + \varepsilon \eta'^\mu] &= \int \diff l\, \mathcal{L}(l, x^\mu + \varepsilon \eta^\mu, x'^\mu + \varepsilon \eta'^\mu)
		\end{align*}
	
	Da angenommen wird, dass $x^\mu$ $S$ maximiert, muss $S_\varepsilon$ maximal werden, wenn $\varepsilon$ null ist, d.h.
		\begin{align*}
			0 &\stackrel{!}{=} \left.\frac{\diff S_\varepsilon}{\diff\varepsilon} \right|_{\varepsilon = 0} = \int \diff l\, \left. \frac{\diff \mathcal{L}(l, x^\mu + \varepsilon\eta^\mu, x'^\mu + \varepsilon \eta'^\mu}{\diff\varepsilon} \right|_{\varepsilon = 0} = \\
			&= \int \diff l\, \lr{ \frac{\partial \mathcal{L}(l, x^\mu + \varepsilon\eta^\mu, x'^\mu + \varepsilon \eta'^\mu)}{\partial (x^\sigma + \varepsilon\eta^\sigma)}\frac{\partial (x^\sigma + \varepsilon\eta^\sigma)}{\partial\varepsilon} + \frac{\partial \mathcal{L}(l, x^\mu + \varepsilon\eta^\mu, x'^\mu + \varepsilon \eta'^\mu)}{\partial (x'^\sigma + \varepsilon\eta'^\sigma)}\frac{\partial (x'^\sigma + \varepsilon\eta'^\sigma)}{\partial\varepsilon} } = \\
			&= \int \diff l\, \lr{ \frac{\partial \mathcal{L}(l, x^\mu, x'^\mu)}{\partial x^\sigma}\eta^\sigma + \frac{\partial \mathcal{L}(l, x^\mu, x'^\mu)}{\partial x'^\sigma}\eta'^\sigma } = \\
			&= \int \diff l\, \lreck{ \frac{\partial \mathcal{L}(l, x^\mu, x'^\mu)}{\partial x^\sigma}\eta^\sigma + \frac{\diff}{\diff l} \lr{ \frac{\partial \mathcal{L}(l, x^\mu, x'^\mu)}{\partial x'^\sigma}\eta^\sigma } - \frac{\diff}{\diff l} \lr{ \frac{\partial \mathcal{L}(l, x^\mu, x'^\mu)}{\partial x'^\sigma} } \eta^\sigma } = \\
			&= \int \diff l\, \lreck{ \lr{ \frac{\partial\mathcal{L}}{\partial x'^\sigma}\eta^\sigma - \frac{\diff}{\diff l} \frac{\partial\mathcal{L}}{\partial x'^\sigma} } \eta^\sigma + \frac{\diff}{\diff l} \lr{ \frac{\partial \mathcal{L}}{\partial x'^\sigma}\eta^\sigma } } = \\
			&= \int \diff l\, \underbrace{ \lr{ \frac{\partial\mathcal{L}}{\partial x'^\sigma} - \frac{\diff}{\diff l} \frac{\partial\mathcal{L}}{\partial x'^\sigma} } }_{\stackrel{!}{=}0} \eta^\sigma + \underbrace{ \left.\frac{\partial \mathcal{L}}{\partial x'^\sigma}\eta^\sigma \right|_{l_i}^{l_f} }_{=0}
		\end{align*}

	Damit diees Integral null wird, muss der Vorfaktor von $\eta^\sigma$ null werden, was den Euler-Lagrange-Gleichungen entspricht.
		$$\Leftrightarrow \frac{\diff}{\diff l} \frac{\partial\mathcal{L}}{\partial x'^\sigma} - \frac{\partial\mathcal{L}}{\partial x'^\sigma} = 0 \qquad\blacksquare$$


\hrule
}


\newpage


\section{Elektromagnetismus}
Die Lagrange-Dichte des elektromagnetischen Feldes ist
	\begin{equation}
		\mathcal{L} = \mathcal{L}_\text{EM} + \mathcal{L}_J = -\frac{1}{4} F_{\mu\nu}F^{\mu\nu} + J^\mu A_\mu
	\end{equation}

mit dem elektromagnetischen Feldstärketensor
	\begin{align}
		F_{\mu\nu} &= 
		\begin{pmatrix}
			0 & E^1 & E^2 & E^3 \\
			-E^1 & 0 & B^3 & -B^2 \\
			-E^2 & -B^3 & 0 & B^1 \\
			-E^3 & B^2 & -B^1 & 0
		\end{pmatrix} = \\
		&= \partial_\mu A_\nu - \partial_\nu A_\mu \\
		F^{\mu\nu} &= 
		\begin{pmatrix}
			0 & -E^1 & -E^2 & -E^3 \\
			E^1 & 0 & -B^3 & B^2 \\
			E^2 & B^3 & 0 & -B^1 \\
			E^3 & -B^2 & B^1 & 0
		\end{pmatrix}
	\end{align}

Die Wirkung $S = \int \diff^4 x \sqrt{-g}\, \mathcal{L}$ wird extremal falls
	\begin{equation}
		\partial_\mu F^{\mu\nu} = -J^\nu
	\end{equation}

gilt \textcolor{gruen}{(folgt aus Variationsrechnung bzw. Euler-Lagrange-Gl.)}, was den Maxwell-Gleichungen
	\begin{align}
		\vec{\nabla} \vec{E} &= \rho \\
		\vec{\nabla} \vec{B} &= 0 \\
		\vec{\nabla} \times \vec{E} &= -\partial_t \vec{B} \\
		\vec{\nabla} \times \vec{B} &= \vec{J} + \partial_t \vec{E}
	\end{align}

entspricht. {\color{gruen}Dabei folgen nur 2 der 4 Gln. aus $\partial_\mu F^{\mu\nu} = -J^\nu$. Die anderen folgen aus einer mathematischen Identität, die als $\partial_\mu \tilde{F}^{\mu\nu} = 0$ geschrieben werden kann. Dabei ist $\tilde{F}$ der duale elektromagnetische Felstärketensor
	\begin{equation}
		\tilde{F}_{\mu\nu} := -\frac{1}{2}\varepsilon_{\mu\nu\lambda\sigma} F^{\lambda\sigma} =
		\begin{pmatrix}
			0 & -B^1 & -B^2 & -B^3 \\
			B^1 & 0 & -E^3 & E^2 \\
			B^2 & E^3 & 0 & -E^1 \\
			B^3 & -E^2 & E^1 & 0
		\end{pmatrix}
	\end{equation}

}


\newpage


\section{Signalaustausch}
Betrachte ein de-Sitter-Universum, d.h. es gelte die de-Sitter-Metrik mit $a(t) = e^{H\cdot t}$.\\
Betrachte die Situation, dass zur Zeit $t=t_S$ und am Ort $x=0$ ein Signal zum Ort $x=R$ geschickt wird, das zu $t=t_U$ zurückgeschickt wird und zu $t=t_R$ wieder bei $x=0$ ankommt.

\begin{figure}[H]
	\centering
	\fbox{
		\includegraphics[width=0.3\textwidth]{Grafiken/ART Signalaustausch.pdf}
	}
	\caption{Skizze Signalaustausch}
\end{figure}



Für einen Lichtstrahl gilt der räumliche Abstand
	\begin{equation}
		R = \frac{1}{H} \lr{ e^{-H t_S} - e^{-H t_U} }
	\end{equation}

Für $R > R_\text{grenz}$ kommt kein Signal mehr an, da $t_U\to\infty$ wird
	\begin{equation}
		R_\text{grenz} = \frac{1}{H} e^{-H t_S}
	\end{equation}

Der gesamte Abstand, also Hin- und Rückweg, ist
	\begin{equation}
		D(R,t_S) = {\color{gruen}\frac{t_R-t_S}{2}} = -\frac{1}{2H} \ln\lr{ 1 - 2HRe^{+H\cdot t_S}}
	\end{equation}

Dieser wird für $R = \frac{1}{2}R_\text{grenz}$ unendlich, d.h. das Signal kommt nicht mehr zurück.














\newpage


\section{Kovariante Ableitung}
Das besondere an der kovarianten Ableitung ist, dass sie, angewendet auf einen Tensor, wieder einen Tensor ergibt. Dies ist bei der partiellen Ableitung nicht der Fall.\\
Die kovariante Ableitung eines kontravarianten Vektores, kovarianten Vektors und eines ($m,n$)-Tensor ist
	\begin{align}
		D_\lambda W^\mu &:= \partial_\lambda W^\mu + \Gamma^\mu_{\lambda\alpha}W^\alpha \equiv W^{\mu}{}_{;\lambda} \\
		D_\lambda W_\mu &:= \partial_\lambda W_\mu - \Gamma^\alpha_{\lambda\mu}W_\alpha \equiv W_{\mu;\lambda} \\
		D_\lambda T^{\mu_1,...,\mu_m}{}_{\nu_1,...,\nu_n} &:= \partial_\lambda T^{\mu_1,...,\mu_m}{}_{\nu_1,...,\nu_n} + \Gamma^{{\color{red}\mu_1}}_{\lambda{\color{gruen}\alpha}} T^{{\color{gruen}\alpha},...,\mu_m}{}_{\nu_1,...,\nu_n} + ... - \\
		 &\phantom{:=} - \Gamma^{{\color{gruen}\alpha}}_{\lambda{\color{blue}\nu_1}} T^{\mu_1,...,\mu_m}{}_{{\color{gruen}\alpha},...,\nu_n} - ... \notag
	\end{align}

\textbf{Rechenregeln}
	\begin{itemize}
		\item \textbf{Produktregel}
			$$(U_\mu W^\mu)_{;\lambda} = U_{\mu;\lambda}W^\mu + U_\mu W^\mu{}_{;\lambda}$$
		
		\item \textbf{Rotation}\\
		Die kovariante Rotation ist die gewöhnliche Rotation
			$$\partial_\mu W_\nu - \partial_\nu W_\mu = D_\mu W_\nu - D_\nu W_\mu$$
	\end{itemize}

Insbesondere gilt
	\begin{align*}
		D_\lambda g_{\mu\nu} &= 0 \\
		D_\lambda g^{\mu\nu} &= 0
	\end{align*}


\newpage


\section{Riemann-Tensor}
Der Riemann-Tensor wie folgt definierte Tensor hat aufgrund seiner Symmetrie- und Zyklizitätseigenschaften statt 256 nur 20 unabhängige Komponenten.
	\begin{align}
		R^\sigma{}_{\rho\mu\nu} &= \partial_\mu \Gamma^\sigma_{\nu\rho} + \Gamma^\sigma_{\mu\kappa}\Gamma^\kappa_{\nu\rho} - \partial_\nu \Gamma^\sigma_{\mu\rho} - \Gamma^\sigma_{\nu\kappa}\Gamma^\kappa_{\mu\rho} \\
		R_{\tau\rho\mu\nu} &:= g_{\tau\lambda} R^\lambda{}_{\rho\mu\nu}
	\end{align}

\begin{itemize}
	\item \textbf{Symmetrien}
		\begin{align}
			R_{\sigma\rho{\color{red}\mu}{\color{blue}\nu}} &= -R_{\sigma\rho{\color{blue}\nu}{\color{red}\mu}} \\
			R_{{\color{red}\sigma}{\color{blue}\rho}\mu\nu} &= -R_{{\color{blue}\rho}{\color{red}\sigma}\mu\nu} \\
			R_{{\color{red}\sigma\rho}{\color{blue}\mu\nu}} &= R_{{\color{blue}\mu\nu}{\color{red}\sigma\rho}}
		\end{align}

		\item \textbf{Zyklizität} in $\rho\mu\nu$
			\begin{equation}
				R_{\sigma\rho\mu\nu} + R_{\sigma\mu\nu\rho} + R_{\sigma\nu\rho\mu} = 0
			\end{equation}
			
		\item \textbf{Komponenten}\\
		Der Riemann-Tensor hat für $D$ Indizes
			$$\frac{D^2(D^2-1)}{12}$$

		Komponenten
\end{itemize}


{\color{violet}
\hrule
	\fbox{Herleitung}\\
	Der Riemann-Tensor ist der Kommutator zweier kovarianter Ableitungen, die auf einen kovarianten Vektor wirken.\\
	Tipp: Führe zuerst die ''äußeren'' Ableitungen aus.
		\begin{align*}
			[D_\mu,D_\nu] S_\rho &= D_\mu D_\nu S_\rho - D_\nu D_\mu S_\rho = \\
			&\stackrel{\comment{D_\mu,D_\nu}}{=} \partial_\mu (D_\nu S_\rho) - \cancel{\Gamma^\alpha_{\mu\nu} D_\alpha S_\rho} - \Gamma^\alpha_{\mu\rho} D_\nu S_\alpha - \lr{ \partial_\nu (D_\mu S_\rho) - \cancel{\Gamma^\alpha_{\nu\mu} D_\alpha S_\rho} - \Gamma^\alpha_{\nu\rho} D_\mu S_\alpha } = \\
			&\stackrel{\comment{D_\nu,D_\mu}}{=} \partial_\mu (\cancel{\partial_\nu S_\rho} - \Gamma^\alpha_{\nu\rho} S_\alpha) - \Gamma^\alpha_{\mu\rho} (\partial_\nu S_\alpha - \Gamma^\beta_{\nu\alpha} S_\beta) - \\
			&\phantom{=} - \partial_\nu (\cancel{\partial_\mu S_\rho} - \Gamma^\alpha_{\mu\rho}S_\alpha) + \Gamma^\alpha_{\nu\rho} (\partial_\mu S_\alpha - \Gamma^\beta_{\mu\alpha} S_\beta) = \\
			&\stackrel{\commenttext{Prod.reg.}}{=} -\partial_\mu (\Gamma^\alpha_{\nu\rho}) S_\alpha - \cancel{\Gamma^\alpha_{\nu\rho} \partial_\mu S_\alpha} - \Gamma^\alpha_{\mu\rho} (\bcancel{\partial_\nu S_\alpha} - \Gamma^\beta_{\nu\alpha} S_\beta) + \\
			&\phantom{=} + \partial_\nu (\Gamma^\alpha_{\mu\rho}) S_\alpha + \bcancel{\Gamma^\alpha_{\mu\rho} \partial_\nu S_\alpha} + \Gamma^\alpha_{\nu\rho} (\cancel{\partial_\mu S_\alpha} - \Gamma^\beta_{\mu\alpha} S_\beta) = \\
			&= -\partial_\mu(\Gamma^\alpha_{\nu\rho})S_\alpha + \Gamma^\alpha_{\mu\rho}\Gamma^\beta_{\nu\alpha} S_\beta + \partial_\nu (\Gamma^\alpha_{\mu\rho}) S_\alpha - \Gamma^\alpha_{\nu\rho}\Gamma^\beta_{\mu\alpha} S_\beta = \\
			&= \lr{ \partial_\nu \Gamma^\alpha_{\mu\rho} + \Gamma^\alpha_{\nu\kappa}\Gamma^\kappa_{\mu\rho} - \partial_\mu\Gamma^\alpha_{\nu\rho} - \Gamma^\alpha_{\mu\kappa}\Gamma^\kappa_{\nu\rho} } S_\alpha =: \\
			&=: (R^\alpha{}_{\rho\mu\nu}) S_\rho
		\end{align*}
\hrule
}




\subsection{Ricci-Tensor}
	\begin{equation}
		R_{\mu\nu} := g^{{\color{gruen}\tau}{\color{gruen}\mu}} R_{{\color{gruen}\tau}\rho{\color{gruen}\mu}\nu} = R^{\color{gruen}\mu}{}_{\rho{\color{gruen}\mu}\nu} = {\color{red}\text{Ergänzen Riemann-Tensor einsetzen}}
	\end{equation}


\subsection{Ricci-Skalar}
	\begin{equation}
		R := g^{\mu\nu}R_{\mu\nu}
	\end{equation}

\vspace{0.5cm}

{\color{blue}
\hrule
	\fbox{Beispiel}\\
	ges: Riemann-Tensor, Ricci-Tensor und Ricci-Skalar $S^2$ Kugel\\
	\\
	Die $S^2$ Kugel hat die Metrik
		$$g_{\mu\nu} =
		\begin{pmatrix}
			1 & 0 \\
			0 & \sin^2\theta
		\end{pmatrix}$$

	Die Christoffel-Symbole sind
		\begin{align*}
			\Gamma^\rho_{\mu\nu} &= \frac{1}{2}g^{\rho\alpha} (\partial_\mu g_{\alpha\nu} + \partial_\nu g_{\mu\alpha} - \partial_\alpha g_{\mu\nu}) \stackrel{\comment{\alpha \stackrel{!}{=} \rho}}{=} \frac{1}{2}g^{\rho\rho} (\partial_\mu g_{\rho\nu} + \partial_\nu g_{\mu\rho} - \partial_\rho g_{\mu\nu}) \\
			\bullet \Gamma^\theta_{\mu\nu} &= \frac{1}{2} \cdot \frac{1}{1} \cdot (\partial_\mu g_{\theta\nu} + \partial_\nu g_{\mu\theta} - \partial_\theta g_{\mu\nu}) \\
			\Rightarrow \Gamma^\theta_{\varphi\varphi} &= \frac{1}{2} (-\partial_\theta g_{\varphi\varphi}) = -\frac{1}{2} \partial_\theta (\sin^2\theta) = -\frac{1}{2} \cdot 2\sin\theta\cos\theta = -\sin\theta\cos\theta \\
			\bullet \Gamma^\varphi_{\mu\nu} &= \frac{1}{2} \cdot \frac{1}{\sin^2\theta} \cdot (\partial_\mu g_{\varphi\nu} + \partial_\nu g_{\mu\varphi} - \partial_\varphi g_{\mu\nu}) \\
			\Rightarrow \Gamma^\varphi_{\theta\varphi} &= \frac{1}{2} \cdot \frac{1}{\sin^2\theta} \cdot \partial_\theta g_{\varphi\varphi} = \frac{1}{2} \cdot \frac{1}{\sin^2\theta} \cdot \partial_\theta \sin^2\theta = \frac{1}{2} \cdot \frac{1}{\sin^2\theta} \cdot 2\sin\theta\cos\theta = \\
			&= \frac{\cos\theta}{\sin\theta} = \cot\theta
		\end{align*}
	
	\dotfill
	
	Da es 2 Indizes gibt, gibt es $\frac{2^2(2^2-1)}{12} = 1$ unabhängige Komponente des Riemann-Tensors.
		\begin{align*}
			R^\mu{}_{\nu\rho\sigma} &= \partial_\rho \Gamma^\mu_{\nu\sigma} + \Gamma^\mu_{\rho\kappa}\Gamma^\kappa_{\nu\sigma} - \partial_\sigma \Gamma^\mu_{\nu\rho} - \Gamma^\mu_{\sigma\kappa}\Gamma^\kappa_{\nu\rho} \\
			\Rightarrow R^\theta{}_{\varphi\theta\varphi} &= \partial_\theta \Gamma^\theta_{\varphi\varphi} + \cancel{\Gamma^\theta_{\theta\kappa}}\Gamma^\kappa_{\varphi\varphi} - \cancel{\partial_\varphi \Gamma} - \Gamma^\theta_{\varphi\kappa}\Gamma^\kappa_{\theta\varphi} = \partial_\theta (-\sin\theta\cos\theta) - (-\sin\theta\cos\theta)\cdot \cot\theta = \\
			&= -(\cos^2\theta - \sin^2\theta) + \cos^2\theta = \sin^2\theta
		\end{align*}
	
	\dotfill
	
	Die Ricci-Tensoren sind somit
		\begin{align*}
			R_{\varphi\varphi} &= R^\theta{}_{\varphi\theta\varphi} = \sin^2\theta \\
			R_{\theta\theta} &= R^\varphi{}_{\theta\varphi\theta} = g^{\varphi\varphi} g_{\theta\theta} R^\theta{}_{\varphi\theta\varphi} = \frac{1}{\sin^2\theta} \cdot 1 \cdot \sin^2\theta = 1 \\
			R_{\theta\varphi} &= R^\alpha{}_{\theta\alpha\varphi} = R^\theta{}_{\theta\theta\varphi} + R^\varphi{}_{\theta\varphi\varphi} = 0 = R_{\varphi\theta}
		\end{align*}
	
	\dotfill
	
	Der Ricci-Skalar ist dann
		\begin{align*}
			R &= g^{\mu\nu}R_{\mu\nu} = g^{\theta\theta}R_{\theta\theta} + g^{\theta\varphi}R_{\theta\varphi} + g^{\varphi\theta}R_{\varphi\theta} + g^{\varphi\varphi}R_{\varphi\varphi} = \\
			&= 1\cdot 1 + \frac{1}{\sin^2\theta} \cdot \sin^2\theta = 2
		\end{align*}
	

\hrule
}





\newpage


\section{Einstein-Gleichungen}
Die Wirkung des Universums wird durch die Einstein-Hilbert- und die Energie-Materie-Wirkung beschrieben
	\begin{align}
		S &= S_\text{EH} + S_\text{EM} = \frac{c^3}{16\pi G} \int \diff^4x\,\sqrt{-g}\, R - \frac{1}{c} \int \diff^4x\,\sqrt{-g}\, \frac{1}{2} T^{\mu\nu} g_{\mu\nu}
	\end{align}

Variation nach der Metrik $g_{\mu\nu}$ ergibt die \textbf{Einstein'schen Feldgleichungen}
	\begin{equation}
		G_{\mu\nu} := R_{\mu\nu} - \frac{1}{2} g_{\mu\nu} R = \frac{8\pi G}{c^4} T_{\mu\nu}
	\end{equation}

mit dem Einstein-Tensor $G_{\mu\nu}$ und dem Energie-Materie-Tensor
	\begin{equation}
		T^{\mu\nu} :=
		\begin{pmatrix}
			\text{Energiedichte} & (\text{Impulsdichte})^2 \\
			\text{Impusldichte} & \text{Impulsfluss}
		\end{pmatrix}
	\end{equation}


Im Vakuum wird die Einstein-Gleichung zu
	\begin{align*}
		R_{\mu\nu} - \frac{1}{2}g_{\mu\nu} R &= 0 \\
		\stackrel{\comment{\cdot g^{\mu\nu}}}{\Leftrightarrow} -R &= 0 \\
		\stackrel{\commenttext{R in erste Gl.}}{\Leftrightarrow} R_{\mu\nu} &= 0
	\end{align*}


\newpage


\section{Schwarze Löcher}
Die Metrik eines nichtrotierenden, statischen Körpers ist die Schwarzschild-Metrik
	\begin{align}
		\diff s^2 &= -\lr{ 1 - \frac{2GM}{c^2r} } \diff t^2 + \frac{1}{1 - \frac{2GM}{c^2r}} \diff r^2 + r^2 \diff\theta^2 + r^2\sin^2\theta \diff\varphi^2 \\
		\stackrel{\commenttext{Koordinatentrafo}}{\Leftrightarrow} \diff s^2 &= -\frac{r - r_S}{r} \lr{ \diff\bar{t} + \diff r } \lr{ \diff\bar{t} - \frac{r+r_S}{r-r_S}\diff r } + r^2\diff\Omega^2
	\end{align}

\textcolor{blue}{Diese folgt durch das lösen der Vakuum-Einstein-Gleichungen mit einer Randbedingung, die den massiven Körper beschreibt.}\\
\\
Für eine Bewegung mit $\diff\Omega = 0$ muss eine der beiden Klammern null werden, d.h.
	\begin{equation}
		\begin{cases}
			\diff \bar{t} &= -\diff r \\
			\diff \bar{t} &= \frac{r+r_S}{r-r_S}\diff r
		\end{cases}
		\Leftrightarrow
		\begin{cases}
			\text{negative Lsg.: immer einlaufend, $\frac{\diff\bar{t}}{\diff r} = -1$ {\color{gruen}(Raum-Zeit-Diagramm)}} \\
			\text{positive Lsg. für $r<r_S$, negative Lsg. für $r>r_S$}
		\end{cases}		
	\end{equation}


\begin{figure}[H]
	\caption{Raum-Zeit-Diagramme eines Lichtstrahls, der sich mit $\diff \Omega = 0$ entweder auf das schwarze Loch zubewegt oder sich davon entfernt; jede Grafik zeigt eine andere Startposition; für $r<r_S$ kann der Strahl nicht mehr entkommen}
	\centering
	\fbox{
		\includegraphics[width=0.75\textwidth]{Grafiken/Lichtkegel Schwarzschild Metrik.pdf}
	}
\end{figure}


\newpage


\section{Kosmologie}
Betrachte das Friedmann-Robertson-Walker-Universum, d.h. die Metrik
	\begin{equation}
		\diff s^2 = -\diff t^2 + a(t)^2 \lr{ \frac{1}{1 - k\frac{r^2}{L^2}} \diff r^2 + \diff\Omega^2}
	\end{equation}

Berechnen der Christoffel-Symbole, Riemann-Tensoren, Ricci-Tensoren und -Skalare ergibt die Einstein'schen Feldgleichungen, unter Berücksichtigung der kosmologischen Konstante
	\begin{equation}
		G_{\mu\nu} =
		\begin{pmatrix}
			3 \frac{\dot{a}^2}{a^2} + 3K\frac{1}{L^2a^2} - \Lambda & 0 \\
			0 & -(2 a\ddot{a} + \dot{a}^2 + \frac{K}{L^2} - a^2\Lambda) \delta_{ij}
		\end{pmatrix}
		=
		8\pi G
		\begin{pmatrix}
			\rho & 0\\
			0 & a^2 p \cdot \delta_{ij} \
		\end{pmatrix}
	\end{equation}

Die 00-Komponente der Tensorgleichung ergibt die erste Friedmann-Gl.
	\begin{equation}
		\frac{\dot{a}^2}{a^2} = \frac{8\pi G}{3} \rho - K\frac{1}{L^2a^2} + \frac{\Lambda}{3}
	\end{equation}

Die $ii$-Komponente liefert die zweite Friedmann-Gl.
	\begin{equation}
		\frac{\ddot{a}}{a} = -\frac{4\pi G}{3}(\rho + 3p) - \frac{\Lambda}{3}
	\end{equation}


\subsection{Urknall}
Aus $\rho, p, a(t) > 0$ folgt, dass die Beschleunigung $a(t) < 0$. Das Universum ist \underline{nicht statisch}.\\
\\
Die experimentelle Beobachtung von Hubble zeigt, dass
	\begin{equation}
		\dot{a} > 0
	\end{equation}

Aus diesen beiden Erkenntnissen folgt, dass $a(t)$ zu einem bestimmten Zeitpunkt 0 gewesen sein muss.


\subsection{Kosmologisches Prinzip}
Das Universum ist auf Längenskalen $\gg$ Galaxiengrößen homogen und isotrop.\\
D.h. es gibt keinen ausgezeichneten Ort und keine ausgezeichnete Richtung.








\newpage


\section{Stichwortverzeichnis}
\begin{xltabular}{\linewidth}{p{6cm}|X}
	\hline\hline
		Begriff & Erklärung \\
	\hline
		Äquivalenzprinzip (schwach) & Besagt, dass die träge Masse in $\vec{F} = m\vec{a}$ gleich der schweren Masse in $F_G = G \frac{mM}{r^2}$ ist \\
	\hline
		Äquivalenzprinzip (stark) & Besagt, dass in einer infinitesimal kleinen Raum-Zeit-Umgebung kein Experiment durchgeführt werden kann, das darüber aufklärt, ob man sich in einem beschleunigten System oder in einem System im Gravitationsfeld befindet.\\
	\hline
		lichtartig &  \\
	\hline
		raumartig &  \\
	\hline
		Relativitätsprinzip &  \\
	\hline
		zeitartig &  \\
	\hline\hline
\end{xltabular}
\end{document}
